% Generated draft: requires mathematical and source review.
\documentclass[11pt]{article}
\usepackage[margin=25mm]{geometry}
\usepackage{fontspec}
\setmainfont{Latin Modern Roman}
\usepackage{amsmath,amssymb,mathtools,unicode-math}
\setmathfont{Latin Modern Math}
\usepackage{enumitem,needspace,xurl,hyperref}
\setlength{\parindent}{0pt}
\setlength{\parskip}{5pt}
\setlength{\emergencystretch}{3em}
\newcommand{\N}{\mathbb{N}}
\newcommand{\Z}{\mathbb{Z}}
\newcommand{\Q}{\mathbb{Q}}
\newcommand{\R}{\mathbb{R}}
\newcommand{\C}{\mathbb{C}}
\newcommand{\F}{\mathbb{F}}
\newcommand{\A}{\mathbb{A}}
\newcommand{\PP}{\mathbb{P}}
\newcommand{\E}{\mathbb{E}}
\newcommand{\eps}{\varepsilon}
\newcommand{\dd}{\,\mathrm{d}}
\begin{document}
\section*{UnsolvedMath ID: 800876}
\textbf{Source problem code:} catalogue entry
\subsection{Definitions and mathematical statement}
The supplied record attributes to Demailly the question
\[
(dd^c\varphi_j)^n(0)\longrightarrow(dd^c\varphi)^n(0)
\qquad (j\to\infty).
\]
Here the approximation sequence \((\varphi_j)\), the function \(\varphi\), their domain, the hypotheses on their singularities, the normalization of \(d^c\), and the meaning and existence of the displayed Monge--Amp\`ere quantities are not defined in the supplied excerpt.

A conditional interpretation of the question can be made precise as follows. Let \(n\geq1\), let \(U\subseteq\C^n\) be an open neighborhood of the origin, and suppose that the intended functions and Monge--Amp\`ere construction yield positive Borel measures
\[
\mu_j=(dd^c\varphi_j)^n,\qquad \mu=(dd^c\varphi)^n
\]
on \(U\), each finite on compact subsets of \(U\). Define the atomic mass of such a measure \(\nu\) at the origin by \(\nu(\{0\})\). If the source's notation \((dd^c\psi)^n(0)\) denotes this atomic mass, the question is whether the intended approximation satisfies
\[
\lim_{j\to\infty}\mu_j(\{0\})=\mu(\{0\}).
\]
This is only a conditional reading: the supplied excerpt does not establish these hypotheses, specify the approximation, identify the Monge--Amp\`ere construction, or confirm the atomic-mass interpretation. Consequently, a fully quantified statement for the intended class of functions cannot be recovered from this record alone.
\subsection{Short English statement}
Demailly asks whether the Monge--Amp\`ere masses at the origin of the intended approximants \(\varphi_j\) converge to that of \(\varphi\). The excerpt omits the approximation construction and analytic hypotheses, and does not define the point-mass notation. The displayed precise version is therefore conditional and requires human mathematical review.
\subsection{Sources}
\begin{itemize}
\item Local UnsolvedMath registry, permanent ID \(800876\), source code \texttt{catalogue entry}, title \emph{Open Problems in Pluripotential Theory --- Question 7}; supplied set: AMR Open Problem Lists. The registry labels this problem \emph{open}; this label is supplied source data, not an independently verified assessment.
\item Dinew, Guedj, Zeriahi, \emph{Open Problems in Pluripotential Theory} (2015), Question 7. Supplied source URL: \url{https://arxiv.org/abs/1511.00705}. The supplied provenance records access on 2026-10-16 and extraction method \texttt{arxiv\_latex}; these are supplied provenance statements, not a claim of independent consultation.
\item The supplied question is explicitly attributed to Demailly. Its introductory sentence attributes to Demailly in reference [Dem93] the observation that the Ohsawa--Takegoshi theorem gives a very accurate approximation of \(\varphi\) by the \(\varphi_j\). The record supplies neither a definition of this approximation nor bibliographic details for [Dem93]; the observation is retained as source attribution without independent verification.
\end{itemize}
\end{document}
